\documentclass[12pt]{article}
\usepackage[utf8]{inputenc}
\usepackage[T1]{fontenc}
\usepackage{hyperref}
\title{Yuracare Language Resources for K-12 Curriculum Development}
\author{Compiled by Claude Sonnet 5.5 (claude-sonnet-5-5) for Andr\'as Kornai}
\date{}
\begin{document}
\maketitle

\section{Introduction}
Yuracare (also spelled Yurakar\'e; ISO 639-3 yuz) is an endangered language isolate of Cochabamba and Beni, Bolivia. Speaker estimates vary widely, and the DoBeS archive notes that no reliable figures exist. This report lists only the resources that appeared in web search results. Little material exists that is designed for schools.

\section{Primary Educational Resources}
No Yuracare primers, textbooks or school programmes turned up in the searches. A Wikipedia article on the Yuracar\'e people says a first Yuracare--Spanish dictionary and grammar appeared in 2025. I found no second source for this, so it is unverified:
\url{https://en.wikipedia.org/wiki/Yuracar%C3%A9_people}

\section{Linguistic and Reference Materials}
\begin{itemize}
\item Wikipedia overview of the language, which mentions the older missionary-era dictionary by Adam (1893) and Crevels' grammar: \url{https://en.wikipedia.org/wiki/Yuracar%C3%A9_language}
\item DoBeS (MPI Nijmegen) documentation project page on the Yurakar\'e language: \url{https://dobes.mpi.nl/projects/yurakare/language/}
\item DoBeS project page: \url{https://dobes.mpi.nl/?p=1399}
\item Leiden University Scholarly Publications item returned by the search, which may be the Crevels grammar or a related thesis (not verified): \url{https://scholarlypublications.universiteitleiden.nl/access/item%3A3217781/download}
\item Endangered Languages Project entry: \url{https://endangeredlanguages.com/node/98725}
\end{itemize}

\section{Religious and Community Resources}
A community crowdfunding project, ``Preserving Yurakar\'e culture'', aimed to finish and publish the first Yurakar\'e dictionary, with vocabulary on medicine, nature and ceremonies: \url{https://www.awesomefoundation.org/bg/projects/209762-preserving-yurakare-culture}

\section{Digital and Audio Resources}
The DoBeS project pages above are the likely route to archived recordings. The search results did not show individual audio items.

\section{Curriculum Development Context}
No evidence of Yuracare in schools was found. A curriculum would have to start from the community and the documentation projects listed above, and from the grammar and dictionary materials, rather than from existing classroom materials.

\end{document}